% Complete Section 7.1, original PDF 92-93 / printed 82-83.
% block: complete-polynomial-forms-heading page:92 type:section
\subsection{{The complete polynomial spaces of differential forms}}\label{sec:complete-polynomial-forms}

% block: finite-element-assembly-definition page:92 type:ordinary-paragraph
{The spaces $\mathcal P_r\Lambda^k(\mathcal T_h)$ and $\mathcal P_r^-\Lambda^k(\mathcal T_h)$ will be constructed through the usual process of \emph{finite element assembly}. In general, a finite element space is determined by the specification of the triangulation $\mathcal T_h$, and for each element $T\in\mathcal T_h$:}
\begin{itemize}
\item{a finite dimensional space of shape functions $V(T)$ on $T$, and}
\item{a set of degrees of freedom (DOFs) for $V(T)$ (i.e., a basis for its dual space), with each DOF associated to a particular face of $T$.}
\end{itemize}

% block: finite-element-assembled-space page:92 type:ordinary-paragraph
{The assembled finite element space is then defined to consist of all functions on $\Omega$ whose restriction to each $T$ belongs to the shape function space $V(T)$ and for which the DOFs are single-valued (in the sense that when two elements share a common face, the corresponding DOFs on the face take the same value).}

% block: complete-polynomial-shape-functions pages:92,93 type:ordinary-paragraph
{In the present case, the shape function spaces will be spaces of differential $k$-forms with polynomial coefficients. For $\mathcal P_r\Lambda^k(\mathcal T_h)$, the choice of shape function space $V(T)$ is an obvious one, namely, the complete polynomial space}
\[
\mathcal P_r\Lambda^k(T)=\left\{\,\sum_{1\leq\sigma_1<\cdots<\sigma_k\leq n}p_\sigma\,dx^\sigma\ \middle|\ p_\sigma\in\mathcal P_r(T)\,\right\}
\]
{of a differential $k$-form with polynomial coefficients of degree at most $r$. Its dimension is easily calculated:}
\begin{equation}\label{eq:complete-polynomial-form-dimension}
\begin{aligned}
\dim\mathcal P_r\Lambda^k(T)&=\binom nk\times\dim\mathcal P_r(T)\\
&=\binom nk\binom{n+r}n=\binom{r+n}{r+k}\binom{r+k}r.
\end{aligned}
\end{equation}

% block: polynomial-de-rham-complex page:93 type:ordinary-paragraph
{Note that $d\mathcal P_r\Lambda^k\subset\mathcal P_{r-1}\Lambda^{k+1}$, i.e., the exterior derivative lowers the polynomial degree at the same time as it raises the form degree. Therefore, for each $r$ we have a polynomial subcomplex of the de Rham complex:}
\[
0\to\mathcal P_r\Lambda^0\xrightarrow{d}\mathcal P_{r-1}\Lambda^1\xrightarrow{d}\cdots\xrightarrow{d}\mathcal P_{r-n}\Lambda^n\to0.
\]
{(We understand $\mathcal P_s\Lambda^k$ to be $0$ if $s<0$.) Moreover, as proved below in Corollary~\ref{cor:polynomial-de-rham-exactness}, this complex is exact, except for the constant functions in $\mathcal P_r\Lambda^0$, which form the kernel of the first $d$. That is, if $\omega\in\mathcal P_s\Lambda^k$ for some $k>0$ and $d\omega=0$, then $\omega=d\mu$ for some $\mu\in\mathcal P_{s+1}\Lambda^{k-1}$. The proof of this fact will involve a very valuable tool called the \emph{Koszul complex}. The same tool will be used to define DOFs and to define an alternative space of shape functions.}
